% Nouveau EGA V (cote 25): archive cote 25, site page 89 (PDF page 90).
% typescript. Draft compared: none.
% AI-transcribed by Claude Opus 5.5 (claude-opus-5-5, adaptive thinking, effort high), Message Batches API, batch ega25b (msgbatch_01A7n5YuXubUPNK1zsepi8RG), 2026-11-13; not yet checked by eye.
\scanpage{25/89}
% NOTE: the top right corner is cut off at the edge of the scan; only the lower part of a number is visible there (probably the page number)
seule conclusion $p_C$ plat, l'hypothèse CM sur X par une simple hypothèse
dimensionelle. Supposons par exemple f une immersion et X irr de dim m,
de sorte que $p_C$ est qu fini , et comme $X_{k\hand{(y)}}$ et $Q(C_{\hand{y}})$ sont irréd de
même dimension, \var{\xout{le premier CM}}{}{typewriter strike-out}{ega25_25_089-v0001} et le deuxième régulier, $p_C$ ne peut être
plat que si $X_{k\hand{(y)}}$ est CM.

Des propriétés plus fines sont les propriétés différentielles de
la projection conique, notamment pour X lisse sur k et $f:X\longrightarrow P$ non rami-%
\var{\inserted[typed]{étudiées au nº12.}}{étudiées au nº12.}{insertion}{ega25_25_089-v0002}
fié, \var{\xout{Dans le cas où les composantes}}{}{typewriter strike-out}{ega25_25_089-v0003} Rappelons que \var{\xout{la fibre générique de}}{}{typewriter strike-out}{ega25_25_089-v0004}
en dehors d'une partie Z de codimension 1 de $Q(C\ )$, le morphisme
$p_{C_{\hand{y}}}$ \var{\inserted[typed]{sur $X(C_{\hand{y}})$}}{sur $X(C_{\hand{y}})$}{insertion}{ega25_25_089-v0005} est \uline{lisse}, Et une analyse plus détaillée (esquissée au nº12) montre
que si \uline{les dimensions des composantes} de X sont $\geq$ m, alors\hand{)}
(ou montrerait, si nous ne la faisons pas) \var{\xout{que}}{}{typewriter strike-out}{ega25_25_089-v0006} en dehors d'une partie $Z'\subset Z$
de $Q(C_{\hand{y}}\ )$ de codimension $\geq 2$, les fibres $p_C^{-1}(\hand{\zeta}) = \underline{X}^{(m)}_{\hand{\zeta}}$ ne peuvent
avoir au pis que des points \var{\xout{doub}}{}{typewriter strike-out}{ega25_25_089-v0007} singuliers ordinaires (au s ens
géométrique), et en fait (si f est une immersion et X est géom irréd),
\var{\xout{juste}}{}{typewriter strike-out}{ega25_25_089-v0008} au plus \uline{un} tel point, ce dernier \var{\inserted[typed]{étant alors}}{étant alors}{insertion}{ega25_25_089-v0009} \var{\xout{\unsure{ayant}}}{}{typewriter strike-out}{ega25_25_089-v0010} nécessairement \var{\xout{\illegible{1}}}{}{typewriter strike-out}{ega25_25_089-v0011} .
rationnel sur $k(\hand{\zeta})$, - ces assertions étant valables tout au moins si
k est de car 0, \uline{ou} à condition\var{\xout{s}}{}{typewriter strike-out}{ega25_25_089-v0012} de remplacer f par un $\phi_n f$ $(n\geq 2)$ comme
au nº 9. Il y a lieu également de donner des propriétés différentielles
de $p_{C_{\hand{y}}}$ dans le cas où dim X $\leqslant$ m , et où par suite $p_{C_{\hand{y}}}$ est défini sur $X_{k\hand{(y)}}$.
Je me borne à indiquer les propriétés suivantes, dont la démonstration
\var{\inserted{doit}}{doit}{insertion}{ega25_25_089-v0013} être facile et \var{\inserted{est}}{est}{insertion}{ega25_25_089-v0014} laissée à Dieudonné:

\uline{Proposition} 14.4. Supposons $f:X\longrightarrow P$ non ramifié, et dim X $\leqslant$ m.
Soit T un sous-préschéma fini de X, Alors :\var{\xout{la restr}}{}{typewriter strike-out}{ega25_25_089-v0015}
\var{\margin{\unsure{X?}}}{\placedmargin{\unsure{X?}}}{marginal note placed}{ega25_25_089-v0016} a) \var{\inserted{Si f est une immersion,}}{Si f est une immersion,}{insertion}{ega25_25_089-v0017} La restriction de $p_C$ à $T_{k\hand{(y)}}$ est \var{\deleted{un monomorphisme,}}{}{deletion}{ega25_25_089-v0018} \var{\inserted{radiciel i.e.\ ``géométriquement injectif''.}}{radiciel i.e.\ ``géométriquement injectif''.}{insertion}{ega25_25_089-v0019} Si de
plus Y est une partie fermée de X de dimension $\leqslant$ m-1, on a
$p_{C_{\hand{y}}}^{-1}(p_{C_{\hand{y}}}(Y_{k\hand{(y)}})) \cap T_{k\hand{(y)}} = \emptyset$ .
b) Si X est lisse en les points de T, alors $p_{C_{\hand{y}}}$ est non ramifié
en tous point \var{\xout{de}}{}{typewriter strike-out}{ega25_25_089-v0020} $T_{k\hand{(y)}}$, \var{\xout{au dessus de tout point de T}}{}{typewriter strike-out}{ega25_25_089-v0021} \var{\inserted{et \unsure{même|aussi} en les pts de $p_{C_y}^{-1}\,p_{C_y}(T_{k(y)})$.}}{et \unsure{même|aussi} en les pts de $p_{C_y}^{-1}\,p_{C_y}(T_{k(y)})$.}{insertion}{ega25_25_089-v0022}
\mspage{\hand{83}}
